\subsection{线性表示}

设 \(G\) 为李群，\(E\) 为完备可赋范空间，\(\pi\) 为 \(G\) 在 \(E\) 上的一个解析线性表示 (§ 1, no. 2)。相伴态射 \(t \mapsto \langle t, \pi \rangle\)（\(\mathcal{F}^{(\infty)}(G)\) 到 \(\mathcal{L}(E)\) 中的）是代数态射 (第 9 节, 命题 33)，且它在 \(L(G)\) 上的限制是 \(L(\pi)\)。因而 \(L(\pi)\) 是 \(L(G)\) 在 \(E\) 上的一个表示 (Chapter I, § 3, Definition 1)。

命题 38. 把 G 看作通过映射 \((g,x)\mapsto \pi (g)x\) 左作用在 E 上。设 \(b\in \mathrm{E}\)，\(\rho (b)\) 为其轨道映射。把 \(\mathbf{T}_b(\mathbf{E})\) 典范地与 E 认同。对所有 \(t\in \mathbf{L}(\mathbf{G})\)

\[
(\mathbf {L} (\pi) t) (b) = \langle t, \rho (b) \rangle = \rho (b) _ {*} t = t * \varepsilon_ {b}.
\]

特别地，由 \(t\) 在 \(\mathbf{E}\) 上定义的向量场就是场 \(b \mapsto (\mathbf{L}(\pi)t)(b)\)。

\[
\begin{array}{r l} (\mathrm{L} (\pi) t) (b) & = \langle t, g \mapsto \pi (g) b \rangle \\ & = \langle t, \mathrm{Id} _ {\mathrm{E}} \circ \rho (b) \rangle \\ & = \langle \rho (b) _ {*} t, \mathrm{Id} _ {\mathrm{E}} \rangle \quad (Diff. \& Anal. Man., R, 13.2.3) \\ & = \rho (b) _ {*} t. \end{array}
\]

最后，\(\rho(b)_{*}t = t*\varepsilon_{b}\) (第 3 节, 命题 14 (ii))。

命题 39. 设 K 的特征为 0。设 G 为李群，E 为有限维向量空间，\(\pi\) 为 G 在 E 上的解析线性表示。设 \(\mathrm{E}_1\)、\(\mathrm{E}_2\) 为 E 的向量子空间，满足 \(\mathrm{E}_2 \subset \mathrm{E}_1\)。集合 \(\mathrm{G}_1\) 由这样的 \(g \in \mathrm{G}\) 组成：\(\pi(g) x \equiv x \pmod{\mathrm{E}_2}\) 对所有 \(x \in \mathrm{E}_1\) 都成立；它是 G 的李子群，且 \(\mathrm{L}(\mathrm{G}_1)\) 是由 \(a \in \mathrm{L}(\mathrm{G})\)（使得 \(\mathrm{L}(\pi)a\) 把 \(\mathrm{E}_1\) 映入 \(\mathrm{E}_2\)）组成的集合。

这由第 8 节的命题 29 与第 10 节的命题 36 得出。

推论 1. 在命题 39 的记号下，由 \(g \in G\)（满足 \(\pi(g)(\mathrm{E}_1) \subset \mathrm{E}_1\)）组成的集合是 \(G\) 的李子群，其李代数是由 \(a \in L(G)\)（使得 \(L(\pi)a\) 把 \(E_1\) 映入 \(E_1\)）组成的集合。

我们对 \(E_{1} = E_{2}\) 的情形应用命题 39。

推论 2. 设 G、E、\(\pi\) 如命题 39。设 F 为 E 的子集。由 \(g \in \mathbf{G}\)（满足 \(\pi(g)x = x\)，对所有 \(x \in \mathbf{F}\)）组成的集合是 G 的李子群，其李代数是由 \(a \in \mathbf{L}(\mathbf{G})\)（满足 \((\mathbf{L}(\pi)a)(x) = 0\)，对所有 \(x \in \mathbf{F}\)）组成的集合。

我们对 \(E_{2}=\{0\}\) 且 \(E_{1}\) 为由 F 生成的 E 的向量子空间的情形应用命题 39。

设 \(\pi_1, \pi_2, \ldots, \pi_n\) 为 \(G\) 的解析线性表示。显然，\(\pi\)（诸 \(\pi_1\) 的直和）(Algebra, Chapter VIII, § 13, no. 1) 是 \(G\) 的解析线性表示，且 \(L(\pi)\) 是 \(L(\pi_1), L(\pi_2), \ldots, L(\pi_n)\) 的直和 (Chapter I, § 3, no. 1)。

命题 40. 设 G 为李群，E 为完备可赋范空间，\(\pi\) 为 G 在 E 上的解析线性表示，F 为 E 的一个在 \(\pi(\mathbf{G})\) 下稳定的闭向量子空间。设 K 的特征为 0，或者 F 是 E 的直和因子。

\begin{enumerate}
\def\labelenumi{(\roman{enumi})}
\item
  子表示 \(\pi_1\) 与商表示 \(\pi_2\)（\(\pi\) 的、由 F 定义的）是解析表示。
\item
  F 在 L(π)(L(G)) 下稳定。
\item
  设 \(\rho_{1}\) 与 \(\rho_{2}\) 为由 F 定义的 \(\mathbf{L}(\pi)\) 的子表示与商表示。则 \(\mathbf{L}(\pi_1) = \rho_1, \mathbf{L}(\pi_2) = \rho_2\)。
\end{enumerate}

设 \(\mathrm{A}\) 为由 \(u\in \mathcal{L}(\mathrm{E})\)（满足 \(u(\mathrm{F})\subset \mathrm{F}\)）组成的集合。则 \(\mathrm{A}\) 是 \(\mathcal{L}(\mathrm{E})\) 的闭向量子空间，且 \(\pi\) 取值于 \(\mathbf{A}\)。由关于 \(\mathbf{K}\) 与 \(\mathbf{F}\) 的假设，映射 \(\pi^{\prime}\colon \mathrm{G}\to \mathrm{A}\)（与 \(\pi\) 有同一图像的映射）是解析的 (Differentiable and Analytic Manifolds, R, 5.8.5)。典范映射 \(\theta_{1}\colon \mathrm{A}\to \mathcal{L}(\mathrm{F})\) 与 \(\theta_{2}\colon \mathrm{A}\to \mathcal{L}(\mathrm{E / F})\) 是连续的、线性的，因而是解析的。这就证明了 (i)。映射 \(\mathrm{T_e}(\pi)\) 与 \(\mathrm{T_e}(\pi ')\) 有同一图像，因而 \(\mathrm{L}(\pi)(\mathrm{L}(\mathrm{G}))\subset \mathrm{A}\)，这证明了 (ii)。我们有

\[
\begin{array}{l} \mathrm{T} _ {e} (\pi_ {1}) = \mathrm{T} _ {e} (\theta_ {1} \circ \pi^ {\prime}) = \theta_ {1} \circ \mathrm{T} _ {e} (\pi^ {\prime}) = \rho_ {1} \\ \mathrm{T} _ {e} (\pi_ {2}) = \mathrm{T} _ {e} (\theta_ {2} \circ \pi^ {\prime}) = \theta_ {2} \circ \mathrm{T} _ {e} (\pi^ {\prime}) = \rho_ {2}. \end{array}
\]

命题 41. 设 G 为李群，\(\pi_1, \pi_2, \ldots, \pi_n\)、\(\pi\) 为 G 在完备可赋范空间 \(E_1, E_2, \ldots, E_n\)、E 上的解析线性表示。设

\[
\left(x _ {1}, x _ {2}, \dots , x _ {n}\right) \mapsto x _ {1} x _ {2} \dots x _ {n}
\]

为 \(E_{1} \times E_{2} \times \cdots \times E_{n}\) 到 E 中的连续多重线性映射。设

\[
\pi (g) \left(x _ {1} x _ {2} \dots x _ {n}\right) = \left(\pi_ {1} (g) x _ {1}\right) \left(\pi_ {2} (g) x _ {2}\right) \dots \left(\pi_ {n} (g) x _ {n}\right)
\]

对所有 \(g \in \mathbf{G}\)，\(x_{1} \in \mathrm{E}_{1}, \ldots, x_{n} \in \mathrm{E}_{n}\) 成立。则

\[
(\mathrm{L} (\pi) a) (x _ {1} x _ {2} \dots x _ {n}) = \sum_ {i = 1} ^ {n} x _ {1} x _ {2} \dots x _ {i - 1} ((\mathrm{L} (\pi_ {i}) a) x _ {i}) x _ {i + 1} \dots x _ {n}
\]

对所有 \(a \in \mathrm{L}(\mathrm{G})\)，\(x_1 \in \mathrm{E}_1, \ldots, x_n \in \mathrm{E}_n\) 成立。

作为例子，我们对 n = 2 作计算。

\[
\begin{array}{l l} (\mathrm{L} (\pi) a) (x _ {1} x _ {2}) & = \langle a, g \mapsto \pi (g) (x _ {1} x _ {2}) \rangle \\ & = \langle a, (g \mapsto \pi_ {1} (g) x _ {1}) (g \mapsto \pi_ {2} (g) x _ {2}) \rangle \\ & = \langle a, g \mapsto \pi_ {1} (g) x _ {1} \rangle . x _ {2} + x _ {1}. \langle a, g \mapsto \pi_ {2} (g) x _ {2} \rangle \\ & = ((\mathrm{L} (\pi_ {1}) a) x _ {1}). x _ {2} + x _ {1}. ((\mathrm{L} (\pi_ {2}) a) x _ {2}) \quad \text {(Proposition 38).} \end{array}
\]

推论 1. 设 G 为李群，\(\mathbf{E}_1,\ldots ,\mathbf{E}_{n + 1}\) 为完备可赋范空间，\(\pi_1,\dots ,\pi_{n + 1}\) 为 G 在 \(\mathbf{E}_1,\dots ,\mathbf{E}_{n + 1}\) 上的解析线性表示。设

\[
\mathrm{E} = \mathscr {L} (\mathrm{E} _ {1}, \dots , \mathrm{E} _ {n}; \mathrm{E} _ {n + 1})
\]

为 \(E_{1} \times \cdots \times E_{n}\) 到 \(E_{n+1}\) 中的连续多重线性映射所成的完备可赋范空间 (General Topology, Chapter X, §3, no. 2)。对所有 \(g \in G\)，设 \(\pi(g)\) 为 E 的由下式定义的自同构

\[
(\pi (g) u) (x _ {1}, \dots , x _ {n}) = \pi_ {n + 1} (g) (u (\pi_ {1} (g) ^ {- 1} x _ {1}, \dots , \pi_ {n} (g) ^ {- 1} x _ {n})).
\]

则 \(\pi\) 是 \(\mathbf{G}\) 在 \(\mathbf{E}\) 上的解析线性表示，且

\[
\begin{array}{l} ((\mathrm{L} (\pi) a) u) (x _ {1}, \ldots , x _ {n}) = - \sum_ {i = 1} ^ {n} u (x _ {1}, \ldots , x _ {i - 1}, (\mathrm{L} (\pi_ {i}) a) x _ {i}, x _ {i + 1}, \ldots , x _ {n}) \\ \qquad + (\mathrm{L} (\pi_ {n + 1}) a) (u (x _ {1}, \ldots , x _ {n})) \end{array}
\]

对所有 \(a \in \mathrm{L}(\mathrm{G})\)，\(u \in \mathrm{E}\)，\(x_1 \in \mathrm{E}_1, \ldots, x_n \in \mathrm{E}_n\) 成立。

每个元素 \((A_{1},\ldots,A_{n+1})\)（\(\mathcal{L}(\mathrm{E}_{1})\times\cdots\times\mathcal{L}(\mathrm{E}_{n+1})\) 的）都由公式定义 E 的一个连续自同态 \(\theta(A_{1},\ldots,A_{n+1})\)

\[
(\theta (A _ {1}, \dots , A _ {n + 1}) u) (x _ {1}, \dots , x _ {n}) = A _ {n + 1} (u (A _ {1} x _ {1}, \dots , A _ {n} x _ {n})).
\]

映射 \(\theta\)（从 \(\mathcal{L}(\mathrm{E}_1) \times \cdots \times \mathcal{L}(\mathrm{E}_{n+1})\) 到 \(\mathcal{L}(\mathrm{E})\) 中）是连续的、多重线性的。于是，对所有 \(g \in \mathrm{G}\)，

\[
\pi (g) = \theta (\pi_ {1} (g ^ {- 1}), \dots , \pi_ {n} (g ^ {- 1}), \pi_ {n + 1} (g))
\]

因而 \(\pi\) 是解析的。我们对映射

\[
\left(x _ {1}, \dots , x _ {n}, u\right) \mapsto u \left(x _ {1}, \dots , x _ {n}\right)
\]

（从 \(\mathbf{E}_1\times \dots \times \mathbf{E}_n\times \mathbf{E}\) 到 \(\mathbf{E}_{n + 1}\) 中）应用命题 41。则

\[
\pi_ {n + 1} (g) \left(u \left(x _ {1}, \dots , x _ {n}\right)\right) = (\pi (g) u) \left(\pi_ {1} (g) x _ {1}, \dots , \pi_ {n} (g) x _ {n}\right)
\]

因而

\[
\begin{array}{l} (\mathrm{L} (\pi_ {n + 1}) a) (u (x _ {1}, \ldots , x _ {n})) = \sum_ {i = 1} ^ {n} u (x _ {1}, \ldots , (\mathrm{L} (\pi_ {i}) a) x _ {i}, \ldots , x _ {n}) \\ \qquad \qquad \qquad + ((\mathrm{L} (\pi) a) u) (x _ {1}, \ldots , x _ {n}). \end{array}
\]

当诸 \(E_{i}\) 都是有限维时，表示 \(\mathrm{L}(\pi)\)（\(\mathrm{L}(\mathrm{G})\) 的表示）由表示 \(\mathrm{L}(\pi_{1}),\ldots,\mathrm{L}(\pi_{n+1})\) 按 Chapter I, §3, Proposition 3 的程序导出。

推论 2. 设 G 为李群，\(\pi\) 为 G 在完备可赋范空间 E 上的解析线性表示。则 \(g \mapsto {}^{t}\pi(g)^{-1}\) 是解析线性表示 \(\rho\)（G 在完备可赋范空间 \(\mathcal{L}(\mathrm{E}, \mathrm{K})\dagger\) 上的），且 \(\mathrm{L}(\rho)a = -{}^{t}(\mathrm{L}(\pi)a)\)（对所有 \(a \in \mathrm{L}(\mathrm{G})\)）。

这是推论 1 的特殊情形。

\(\rho\) 称为 \(\pi\) 的逆步表示。

当 E 是有限维时，\(\mathbf{L}(\rho)\) 是 Chapter I, § 3, no. 3 意义下 \(\mathbf{L}(\pi)\) 的对偶表示。

推论 3. 设 G 为李群，\(\pi_1, \ldots, \pi_n\) 为 G 在有限维向量空间 \(E_{1},\ldots,E_{n}\) 上的解析线性表示。则 G 的表示 \(\pi_{1}\otimes\cdots\otimes\pi_{n}\) (Appendix) 是解析的，且 \(\mathrm{L}(\pi_{1}\otimes\cdots\otimes\pi_{n})\) 是 \(\mathrm{L}(\pi_{1}),\ldots,\mathrm{L}(\pi_{n})\) 的张量积。

映射 \((A_{1},\ldots,A_{n})\mapsto A_{1}\otimes\cdots\otimes A_{n}\)（从 \(\mathcal{L}(\mathrm{E}_{1})\times\cdots\times\mathcal{L}(\mathrm{E}_{n})\) 到 \(\mathcal{L}(\mathrm{E}_{1}\otimes\cdots\otimes\mathrm{E}_{n})\) 中）是多线性的，由此可得 \(\pi\) 是解析的。考虑映射 \((x_{1},\ldots,x_{n})\mapsto x_{1}\otimes\cdots\otimes x_{n}\)，它把 \(E_{1}\times\cdots\times E_{n}\) 映到

\[
\mathrm{E} _ {1} \otimes \dots \otimes \mathrm{E} _ {n}.
\]

由命题 41，我们看到

\[
(\mathrm{L} (\pi) a) (x _ {1} \otimes \dots \otimes x _ {n}) = \sum_ {i = 1} ^ {n} x _ {1} \otimes \dots \otimes (\mathrm{L} (\pi_ {i}) a) x _ {i} \otimes \dots \otimes x _ {n}
\]

对所有 \(a \in \mathrm{L}(\mathrm{G})\) 与 \(x_{i} \in \mathrm{E}_{i}\)（\(1 \leqslant i \leqslant n\)）成立。因而 \(\mathrm{L}(\pi)\) 是诸 \(\mathrm{L}(\pi_{i})\) 的张量积。

推论 4. 设 G 为李群，\(\pi\) 为 G 在有限维向量空间 E 上的解析线性表示。则 G 的表示 \(T^{n}(\pi)\)、\(S^{n}(\pi)\) 与 \(\wedge^{n}(\pi)\) (Appendix) 都是解析的，且

\[
\mathrm{L} \left(\mathrm{T} ^ {n} (\pi)\right) = \mathrm{T} ^ {n} (\mathrm{L} (\pi)), \quad \mathrm{L} \left(\mathrm{S} ^ {n} (\pi)\right) = \mathrm{S} ^ {n} (\mathrm{L} (\pi)), \quad \mathrm{L} \left(\wedge^ {n} (\pi)\right) = \wedge^ {n} (\mathrm{L} (\pi)).
\]

这由推论 3 与命题 40 得出。

推论 5. 设 A 为有限维代数。设 K 的特征为 0。A 的自同构群 \(\operatorname{Aut}(\mathbf{A})\) 是 \(\mathbf{GL}(\mathbf{A})\) 的李子群，且 \(\mathbf{L}(\operatorname{Aut}(\mathbf{A}))\) 是 A 的导子李代数。

这由推论 1（取 \(\mathbf{E} = \mathcal{L}(\mathbf{A},\mathbf{A};\mathbf{A})\)）与命题 39 的推论 2（取 \(\mathbf{E}\) 的仅由 \(\mathbf{A}\) 上的乘法组成的子集）得出。

注。我们对 \(\mathrm{G} = \mathbf{GL}(\mathrm{F})\)（F 为完备可赋范空间）、\(\pi_1 = \pi_2 = \mathrm{Id}_{G}\) 且 \(\pi_3\) 为 G 在 K 上的平凡表示的情形应用推论 1。我们得到一个解析表示 \(\pi\)（\(\mathbf{GL}(\mathrm{F})\) 在 \(\mathcal{L}(\mathrm{F},\mathrm{F};\mathrm{K})\) 上的）。设 F 是有限维的且 K 的特征为 0。把命题 39 的推论 2 应用于 \(\pi\)，我们重新得到命题 37 的推论 1 的一部分。

命题 42. 设 G 为李群，X 为解析流形，\((g,x)\mapsto gx\)（相应地 \(xg\)）为 G 在 X 上的解析左（相应地右）作用律，\(x_0\) 为 X 的一个在 G 下不变的点。对所有 \(g\in G\)，设 \(\tau (g)\) 为 X 的自同构 \(x\mapsto gx\)（相应地 \(xg\)），\(\pi (g)\) 为 \(\mathrm{T}_{x_0}(\mathbf{X})\) 的、在 \(x_0\) 处与 \(\tau (g)\) 相切的自同构。

\begin{enumerate}
\def\labelenumi{(\roman{enumi})}
\item
  \(\pi\) 是 G（相应地 \(\mathbf{G}^{\vee}\)）在 \(\mathrm{T}_{x_0}(\mathbf{X})\) 上的解析表示。
\item
  对所有 \(a \in \mathbf{L}(\mathbf{G})\) 与所有 \(\xi_0 \in \mathrm{T}_{x_0}(\mathbf{X}), \mathrm{L}(\pi)a\) . \(\xi_0\) 可按下述方式计算：设 \(\mathbf{D}_a\) 为由 \(a\) 在 \(\mathbf{X}\) 上定义的向量场，\(\xi\) 为 \(\mathbf{C}^1\) 类向量场，定义在 \(x_0\) 的某个开邻域上，满足 \(\xi(x_0) = \xi_0\)；则
\end{enumerate}

\[
\mathrm{L} (\pi) a. \xi_ {0} = - [ \mathrm{D} _ {a}, \xi ] (x _ {0}).
\]

\(\tau(gg') = \tau(g)\tau(g')\)（相应地 \(\tau(g')\tau(g)\)），因而 \(\pi(gg') = \pi(g)\pi(g')\)（相应地 \(\pi(g')\pi(g)\)）。另一方面，由于 TX 是 \(\mathbf{C}^{\omega}\) 类的向量 G-丛 (§ 1, 第 8 节, 命题 16)，\(\pi\) 是解析的，由此得到 (i)。

为证明 (ii)，设 \(G\) 左作用。存在开邻域 \(I\)（\(0\) 在 \(K\) 中的）与解析映射 \(\gamma\)（从 \(I\) 到 \(G\) 中），使得 \(\gamma(0) = e\)，\(T_0(\gamma)1 = a\)。于是 \(D_a\) 是 \(X\) 上的向量场，由映射 \(\phi: (\lambda, x) \mapsto \gamma(\lambda)x\)（从 \(I \times X\) 到 \(X\) 中）所定义 (§ 2, 第 2 节)。若用 \(\phi_{\lambda}\) 表示双射 \(x \mapsto \gamma(\lambda)x\)（从 \(X\) 到 \(X\) 中），则

\[
\begin{array}{r l} \left[ D _ {a}, \xi \right] (x _ {0}) & = \left(\frac {d}{d \lambda} \left(T _ {\phi_ {\lambda} (x _ {0})} (\phi_ {\lambda} ^ {- 1}) \xi (\phi_ {\lambda} (x _ {0}))\right)\right) _ {\lambda = 0} \quad \text {(Diff. \& Anal. Man., R, 8.4.5)} \\ & = \left(\frac {d}{d \lambda} \left(T _ {x _ {0}} (\phi_ {\lambda} ^ {- 1}) \xi_ {0}\right)\right) _ {\lambda = 0} \\ & = \left(\frac {d}{d \lambda} \left(\pi (\gamma (\lambda)) ^ {- 1} \xi_ {0}\right)\right) _ {\lambda = 0}. \end{array}
\]

由于映射 \(\lambda \mapsto \gamma (\lambda)^{-1}\) 与 \(\lambda \mapsto \gamma (-\lambda)\) 在 0 处相切，上式也等于

\[
\begin{array}{r l} & {- \left(\frac {d}{d \lambda} (\pi (\gamma (\lambda)) \xi_ {0})\right) _ {\lambda = 0}} \\ {=} & {- \left(\frac {d}{d \lambda} (\pi \circ \gamma) (\lambda)\right) _ {\lambda = 0} \xi_ {0}} \\ {=} & {- \mathrm{L} (\pi) a. \xi_ {0}.} \end{array}
\]
% semantic-fragments: legacy-input-seam
\relax{}
